% Generated by tools/edn_latex.py from reports/edn.md.
% Do not edit: change reports/edn.md and run `make edn`.
\ednentry{
  id = {EDN-43},
  title = {The persisted contract carries a categorical's level list, and refuses to infer one},
  date = {2026-09-30},
  milestone = {M2: Reproducibility},
  topic = {Data Validation, Feature Encoding},
  participants = {@lukas2510},
  decision = {\texttt{Column} gains a \texttt{levels} field that a \texttt{category} cannot be declared without and that no other dtype may carry, so the artefact the \texttt{features} stage writes is lossless for the only data-dependent dtype it produces. \texttt{Schema.\allowbreak{}validate} compares the level list, not only \texttt{str(dtype)}. \texttt{Schema.\allowbreak{}conform} applies the declared levels instead of inferring them from the frame in front of it, and \texttt{build\_\allowbreak{}features.\allowbreak{}FeatureSpace} is the single public way for a consumer to load the contract and the vocabulary and cast a frame against them.},
  rationale = {\emph{Alternatives considered:}
\begin{itemize}[leftmargin=*, topsep=2pt]
\item \textbf{Option A (chosen): put the levels on \texttt{Column}, in the artefact, and make the cast apply them.}
\newline Pros: the contract becomes self-sufficient -- a consumer holding only the loaded \texttt{Schema} and a frame has a correct cast, which is what \texttt{train} (\#37) and \texttt{evaluate} (\#39) both need. The invariant is structural rather than remembered: \texttt{Column} refuses a \texttt{category} without levels, so no code path can produce a contract that says \texttt{\textquotesingle{}category\textquotesingle{}} and leaves the levels to be guessed. \texttt{validate} now catches a reordered level set, which changes every code in the column while the dtype string stays the same.
\newline Cons: the contract file grows by the level lists, including 279 trims, so the extended set's artefact is larger and noisier to read. Nullable columns keep no \texttt{levels} key, which keeps the noise to the 15 columns that need it.
\item \textbf{Option B: leave the levels in the vocabulary and make the private applier public.}
\newline Pros: no change to the contract at all, and the vocabulary already holds them.
\newline Cons: two objects have to be loaded and kept in step for a cast that is conceptually one contract's business, and \texttt{validate} still cannot check the levels because the schema still does not know them. It also leaves the trap in place: \texttt{conform} would keep silently inferring levels for anyone who used it.
\item \textbf{Option C: document that \texttt{.\allowbreak{}cat.\allowbreak{}codes} is unsafe and leave the mechanism alone.}
\newline Pros: no code change; the pipeline has no active bug today, because pyarrow preserves the dictionary order through Parquet and LightGBM remaps categories by value.
\newline Cons: relies on every future consumer reading the warning. \#37's \texttt{ridge} variant encodes the categoricals itself and \#39's masking sweep rebuilds frames per masked field, and \texttt{.\allowbreak{}cat.\allowbreak{}codes} is the obvious idiom for both, so the trap would be walked into twice within the milestone, each time producing a model scored on a different numbering with no error anywhere.
\end{itemize}
\par
\emph{Why:} The absence was reproduced three ways before it was fixed: \texttt{validate} accepted a frame with reversed levels, \texttt{conform} turned a 50-row subset of 25 levels into 5, and a one-row serving frame for \texttt{make=\allowbreak{}\textquotedbl{}BMW\textquotedbl{}} came back with code 0 where the training code is 3. None of it broke anything today, which is exactly why it was worth fixing now rather than after \#37 and \#39 had each written their own cast. Option B was the smaller change and was rejected because it leaves \texttt{conform} -- the function every stage already calls -- silently wrong; a contract that cannot check the only data-dependent part of itself is not the boundary this module claims to be.},
  ai = {Information seeking, Alternative assessment, Solution generation},
  response = {Accepted with modifications},
  assessment = {The lossy artefact was AI's own design in the \texttt{features} ticket, and an adversarial review by AI found it, reproduced all three consequences by execution and named the two tickets that would walk into it. The modification is the shape of the fix: the review suggested either putting the levels in the contract or having the schema take them from the vocabulary, and the first was chosen because it is the only one that also lets \texttt{validate} check them. The \texttt{Column} invariant and the single \texttt{FeatureSpace} entry point were added on top, so that the property holds by construction rather than by review.},
  aievidence = {Claude Code session on 2026-09-30, adversarial review of \href{https://github.com/mlops-2627q1-mds-upc/MLOPS_Recommenditos/pull/55}{PR \#55}: the review reported that ``\texttt{Schema.\allowbreak{}conform} INFERS a categorical's levels from the data'' and ``\texttt{Vocabulary} does hold the levels, but the only function that applies them is private, so a consumer holding the loaded \texttt{Schema} and a frame has no correct way to cast'', with a reproduction of each case.},
  evidence = {\href{https://github.com/mlops-2627q1-mds-upc/MLOPS_Recommenditos/blob/main/recommenditos/schema.py}{\texttt{recommenditos/\allowbreak{}schema.\allowbreak{}py}}; \href{https://github.com/mlops-2627q1-mds-upc/MLOPS_Recommenditos/blob/main/tests/test_schema.py}{\texttt{tests/\allowbreak{}test\_\allowbreak{}schema.\allowbreak{}py}}; \hyperref[EDN-02]{EDN-02}; \hyperref[EDN-31]{EDN-31}; \href{https://github.com/mlops-2627q1-mds-upc/MLOPS_Recommenditos/issues/37}{issue \#37}; \href{https://github.com/mlops-2627q1-mds-upc/MLOPS_Recommenditos/issues/39}{issue \#39}; \href{https://github.com/mlops-2627q1-mds-upc/MLOPS_Recommenditos/pull/55}{PR \#55}.},
}
